\documentclass[12pt,spanish,letterpaper,oneside]{article}
\def\documenttitle {Estado, representaci\'on y l\'imites al poder}
\def\documentsubtitle {Concepciones del Estado y protecci\'on de la libertad}
\def\documentsubject {Concepciones del Estado y Estado de Derecho}
\def\documentauthor {Martin Jonathan de la Cruz}
\def\coursename {Teor\'ia del Estado y Constituci\'on}
\def\coursecode {020}
\def\universityname {Universidad Abierta y a Distancia de M\'exico}
\def\universityfaculty {Licenciatura en Derecho}
\def\universitydepartment {Teor\'ia del Estado y Constituci\'on}
\def\universitydepartmentimage {departamentos/UnADM}
\def\universitydepartmentimagecfg {height=1.57cm}
\def\universitylocation {Ciudad de M\'exico}
\def\authortable {\begin{tabular}{ll}
Alumno: & Martin Jonathan de la Cruz \\
Matricula: & ES2611202040 \\
Figura docente: & Ana Lilian Sanchez Chavez \\
Semestre/Bloque: & 2 / 2 \\
Semana: & 4 \\
Fecha: & 23 de octubre de 2026
\end{tabular}}
\usepackage{iftex}
\ifLuaTeX
\let\pdfpagewidth\pagewidth
\let\pdfpageheight\pageheight
\fi
\input{template}
\usepackage{fontspec}
\setmainfont{arial.ttf}[Path=C:/Windows/Fonts/,BoldFont=arialbd.ttf,ItalicFont=ariali.ttf,BoldItalicFont=arialbi.ttf]
\setsansfont{arial.ttf}[Path=C:/Windows/Fonts/,BoldFont=arialbd.ttf,ItalicFont=ariali.ttf,BoldItalicFont=arialbi.ttf]
\renewcommand{\familydefault}{\sfdefault}
\setcitestyle{authoryear,open={(},close={)}}
\usepackage[most]{tcolorbox}
\usepackage{fvextra}
\usepackage{attachfile2}
\usepackage{tikz}
\usepackage{eso-pic}
\definecolor{unadmGreen}{HTML}{5F8F3A}
\definecolor{unadmGreenDark}{HTML}{174A3A}
\definecolor{unadmPaper}{HTML}{F6F7F2}
\newtcolorbox{forobox}[1][]{enhanced,breakable,colback=unadmPaper,
colframe=unadmGreen,boxrule=0pt,leftrule=3pt,arc=1.2mm,
left=4mm,right=4mm,top=3mm,bottom=3mm,fonttitle=\bfseries,
coltitle=white,colbacktitle=unadmGreenDark,
attach boxed title to top left={xshift=4mm,yshift=-2mm},
boxed title style={colframe=unadmGreenDark,arc=0.6mm},#1}
\newcommand{\foroTextButton}[1]{\textattachfile[appearance={},description={Texto del foro},mimetype=text/plain]{#1}{\colorbox{unadmGreen}{\small\textcolor{white}{Abrir texto}}}}
\setlength{\emergencystretch}{3em}
% Foro: preguntas guia, participacion completa, lectura, replica real y cierre.
% TXT preparados localmente: no se identifican como mensajes publicados.
\begin{document}
\AddToShipoutPictureBG*{\begin{tikzpicture}[remember picture,overlay]
\node[opacity=0.16,inner sep=0pt] at (current page.center)
{\includegraphics[width=0.72\paperwidth]{img/departamentos/UnADM.pdf}};
\end{tikzpicture}}
\templatePortrait
\templatePagecfg
\onehalfspacing
\begin{abstractd}
El Estado organiza la convivencia, pero su capacidad de imponer orden no
demuestra por s\'i sola representaci\'on social ni legitimidad. El contraste
entre Maquiavelo, Hobbes, Locke, Rousseau y Montesquieu permite distinguir
autonom\'ia de la pol\'itica, seguridad, consentimiento, soberan\'ia y control
del poder. La distribuci\'on de funciones es relevante para proteger la
libertad, aunque requiere legalidad y garant\'ias efectivas. En el orden
constitucional mexicano, soberan\'ia popular, derechos humanos y divisi\'on
de poderes ofrecen criterios concurrentes para valorar la actuaci\'on estatal.
\end{abstractd}
\templateIndex
\templateFinalcfg
\newgeometry{margin=2.5cm}
\fontsize{12}{14.5}\selectfont
\onehalfspacing
\section{Introducci\'on}
El Estado puede producir orden sin representar adecuadamente a la sociedad
y puede invocar representaci\'on sin respetar los l\'imites del poder. Esa
tensi\'on exige distinguir la capacidad de organizar la convivencia de las
razones que legitiman el gobierno y las instituciones que controlan sus
decisiones. El Estado de Derecho no se reduce a la existencia de normas:
implica la sujeci\'on del poder estatal al Derecho \citep[p.~21]{ricciS4}.
Las concepciones de los cinco autores permiten examinar respuestas
hist\'oricamente distintas, no una doctrina uniforme del Estado moderno.

\section{Orden social, soberan\'ia y control constitucional}
\subsection{Concepciones del poder y fundamento del Estado}
La autonom\'ia del pensamiento pol\'itico frente a la teolog\'ia permite
estudiar el gobierno desde sus condiciones humanas e institucionales.
Andrade sit\'ua a Maquiavelo en ese cambio y vincula la reflexi\'on sobre
el Estado con la necesidad de justificar racionalmente la autoridad
\citep[p.~152]{andradeS4}. Esa autonom\'ia no debe convertirse en una
justificaci\'on autom\'atica de cualquier medio empleado para conservar el poder.

Hobbes explica el pacto y la autoridad soberana desde la inseguridad;
Locke rechaza que la protecci\'on de las personas requiera un poder absoluto
e ilimitado. Ambos emplean argumentos contractualistas, pero difieren
en las consecuencias pol\'iticas del consentimiento
\citep[pp.~161--164]{andradeS4}. La diferencia importa: una explicaci\'on
del origen humano del poder no equivale por s\'i sola a gobierno democr\'atico.

Rousseau distingue soberan\'ia y voluntad general de la voluntad particular.
La soberan\'ia no se reduce al mandato ocasional de un gobernante ni a
una suma de intereses privados \citep[libro II, caps.~1--2]{rousseauS4}.
Montesquieu, por su parte, examina el riesgo de concentrar funciones
legislativas, ejecutivas y judiciales y lo relaciona con la libertad
pol\'itica \citep[libro XI, cap.~6]{montesquieuS4}. Sus aportaciones
responden a problemas relacionados, pero no intercambiables: fundamento
del poder, legitimidad del gobierno y prevenci\'on del abuso.

La pregunta sobre los enciclopedistas se interpreta dentro del debate
ilustrado, sin convertir a todos los pensadores mencionados en integrantes
del equipo editor de la Enciclopedia. La preferencia por Montesquieu
se justifica por un criterio expl\'icito: el dise\~no de l\'imites institucionales.
No se afirma que sea el \'unico fundamento del Estado de Derecho.

\subsection{Preguntas sobre orden, representaci\'on y libertad}
\begin{enumerate}
\item \textbf{\textquestiondown El Estado establece un orden en la sociedad o es una representaci\'on de ella?}
Ambas dimensiones pueden relacionarse, pero deben evaluarse por separado:
orden institucional no garantiza participaci\'on ni representaci\'on efectiva.
\item \textbf{\textquestiondown Cu\'al fue la aportaci\'on principal de Maquiavelo, Hobbes, Locke, Rousseau y Montesquieu?}
Autonom\'ia del an\'alisis pol\'itico, seguridad mediante autoridad soberana,
consentimiento y l\'imites, voluntad general y distribuci\'on del poder,
respectivamente; no son una teor\'ia \'unica \citep{andradeS4,rousseauS4,montesquieuS4}.
\item \textbf{\textquestiondown Cu\'al de los enciclopedistas realiz\'o una mayor aportaci\'on al Estado de Derecho?}
Desde el criterio de prevenir la concentraci\'on del poder, Montesquieu
ofrece una aportaci\'on especialmente relevante; necesita complementarse
con legalidad, derechos y garant\'ias \citep[p.~21]{ricciS4}.
\end{enumerate}

\subsection{Estado, representaci\'on y l\'imites al poder}
\begin{forobox}[title={Orden institucional y legitimidad}]
\hfill\foroTextButton{UnADM/licenciatura-en-derecho-unadm/teoria-del-estado-y-constitucion-lde/foro-participacion-Actividad-4.txt}\par\medskip
\VerbatimInput[breaklines=true,breakanywhere=true,fontsize=\footnotesize,baselinestretch=1,breakindent=1.27cm,breaksymbolleft={},breaksymbolright={}]{foro-participacion-Actividad-4.txt}
\end{forobox}
La distinci\'on evita identificar obediencia y consentimiento. Los art\'iculos
39 y 49 permiten situar soberan\'ia y distribuci\'on de competencias; el
art\'iculo 1 agrega deberes de protecci\'on a todas las autoridades
\citep[arts.~1, 39 y 49]{constitucionS4}. Una decisi\'on debe valorarse
tanto por su origen institucional como por los derechos que afecta.

\subsection{Representaci\'on social y controles efectivos}
La intervenci\'on consultada sostiene que orden y representaci\'on van
de la mano y destaca la divisi\'on de poderes. La respuesta distingue
esas funciones y propone evaluar los controles de las actuaciones,
sin identificar a la interlocutora en el encuadre anal\'itico.
\begin{forobox}[title={Legalidad y protecci\'on frente al abuso}]
\hfill\foroTextButton{UnADM/licenciatura-en-derecho-unadm/teoria-del-estado-y-constitucion-lde/foro-replica-Actividad-4.txt}\par\medskip
\VerbatimInput[breaklines=true,breakanywhere=true,fontsize=\footnotesize,baselinestretch=1,breakindent=1.27cm,breaksymbolleft={},breaksymbolright={}]{foro-replica-Actividad-4.txt}
\end{forobox}
La retroalimentaci\'on no niega la relaci\'on entre sociedad y Estado:
precisa que esa relaci\'on puede fallar en la pr\'actica. Un control
efectivo requiere identificar autoridad, norma, derecho afectado y
consecuencia jur\'idica, en vez de considerar suficiente la existencia
formal de poderes separados \citep[p.~21]{ricciS4}.

\clearpage
\section{Conclusiones}
Las concepciones estudiadas distinguen problemas de seguridad, consentimiento,
soberan\'ia, libertad y organizaci\'on del gobierno. Su aporte no consiste
en justificar un poder ilimitado ni en establecer una secuencia de autores
con respuestas id\'enticas. La diferencia entre Hobbes y Locke y la
distinci\'on entre voluntad general y voluntad particular muestran que
el fundamento del Estado admite desacuerdos sustantivos.

Considero que Montesquieu ofrece una aportaci\'on decisiva al control del
poder, siempre que la distribuci\'on de funciones se examine junto con
legalidad y derechos. El Estado debe organizar la convivencia y responder
a la sociedad, pero ambas exigencias necesitan controles verificables.
La representaci\'on no se presume por el mero ejercicio de autoridad y
la libertad no queda asegurada con una divisi\'on exclusivamente formal.
\footnote{GitHub Copilot apoy\'o la organizaci\'on conceptual y el contraste
documental. La postura personal requiere revisi\'on del estudiante.}

\clearpage
\begin{thebibliography}{5}
\bibitem[Andrade S\'anchez(2003)]{andradeS4}
Andrade S\'anchez, E. (2003). \emph{Teor\'ia general del Estado} (2.\textsuperscript{a} ed.). Oxford University Press.
\bibitem[C\'amara de Diputados(2026)]{constitucionS4}
C\'amara de Diputados del H. Congreso de la Uni\'on. (2026). \emph{Constituci\'on Pol\'itica de los Estados Unidos Mexicanos} (\'ultima reforma DOF 7 de octubre de 2026). \url{https://www.diputados.gob.mx/LeyesBiblio/pdf/CPEUM.pdf}
\bibitem[Garc\'ia Ricci(2011)]{ricciS4}
Garc\'ia Ricci, D. (2011). \emph{Estado de derecho y principio de legalidad}. Comisi\'on Nacional de los Derechos Humanos.
\bibitem[Montesquieu(s.~f.)]{montesquieuS4}
Montesquieu. (s.~f.). \emph{Del esp\'iritu de las leyes}.
\bibitem[Rousseau(1999)]{rousseauS4}
Rousseau, J.-J. (1999). \emph{El contrato social o principios de derecho pol\'itico}. elaleph.com.
\end{thebibliography}
\end{document}